% Selección del idioma principal del texto
\ifspanish
	\selectlanguage{spanish}
\else
	\selectlanguage{english}
\fi

% EDITAR: Según el gusto del usuario

% -------------------------
%
% AGRADECIMIENTOS (recomendable máx. 1 pág.)
%
% -------------------------
\phantomsection % Necesario con hyperref
\chapter*{Agradecimientos} % Opción con * para que no aparezca en TOC ni numerada
\addcontentsline{toc}{chapter}{Agradecimientos} % Añade al TOC.

Doce años dan para mucho, y hay gente sin la que nada de esto habría sido posible. En ese tiempo pasé por Affirm - Returnly, Sonnen - Shell, Rollbox - Personio, Bizneo y, actualmente, TeamSystem - Quipu. Acumulé experiencia, responsabilidades, muchos viajes y, sobre todo, compañeros y compañeras que dejaron huella y que, de un modo u otro, también forman parte de lo que hoy entrego. Este trabajo quedó pendiente, por todos estos momentos, es una cuenta sin saldar que hoy, por fin, se cierra.

A mis directores, \textbf{José Antonio Cruz Lemus} y \textbf{Francisco Pascual Romero Chicharro}: gracias por una paciencia que merece figurar en los libros. Desaparecí durante más de un año por cambios de trabajo y vicisitudes varias, y aun así bastaba con un correo para que estuvierais al otro lado, dispuestos a retomar el hilo. Me siento muy afortunado de haberlos tenido como tutores.

A \textbf{Raquel}, mi novia y compañera en este viaje, por las incontables horas leyendo y editando este documento. Sin ella esto seguiría siendo un borrador perdido en alguna carpeta. Gracias por hacerme siempre las preguntas correctas, que son las más difíciles y las más necesarias. Y gracias, sobre todo, por creer en que esto llegaría a buen puerto incluso en los momentos en que yo mismo lo dudaba.

A Prado, mi \textbf{madre}, por insistir tanto en que terminara de una vez. Tenías razón, era una espina que había que quitarse. Espero que ahora puedas dormir un poco más tranquila. Tampoco es casualidad que este trabajo trate sobre la CIE-10, tú, que eres codificadora clínica, has sido mi beta tester favorita y la primera en validar que esto tenía sentido más allá del papel.

A mis \textbf{amigos}, que llevan años con el mismo chascarrillo: \emph{<<¿Qué pasará antes, que Jordi Hurtado se retire o tú acabes la carrera?>>}. Pues bien, señores, la respuesta ya está aquí.

Por último, a todos los \textbf{profesores de la Escuela Superior de Informática de Ciudad Real}. Sin vuestras enseñanzas no habría podido llegar a donde estoy hoy. Fue una etapa genial, y es momento de cerrarla como se merece.

\makeatletter
\begin{flushright}
	\vspace{1.5cm}
	\textit{\@autor}\\
	\@cityTF, \@yearTF
\end{flushright}
\makeatother